\chapter{Tema d'esame del 16 Giugno 2026}
\section*{Esercizio 1}
\textbf{Traccia:}
Si traducano in logica del prim'ordine con uguaglianza le seguenti affermazioni in linguaggio naturale, utilizzando solo i simboli costanti presenti nelle frasi e i simboli di costante e di predicato indicati di seguito. Simboli di costante: 1 e 2. Simboli di predicato: $N(\cdot)="\cdot \text{ è un numero}"$; $Pr(\cdot)="\cdot \text{ è un numero primo}"$; $Succ(\cdot,\circ)="\cdot \text{ è successore di } \circ"$; $Div(\cdot,\circ)="\cdot \text{ è divisibile per } \circ"$; $(\cdot < \circ)="\cdot \text{ minore di } \circ"$.
\begin{enumerate}
    \item[(a)] I numeri primi non formano un insieme denso.
    \item[(b)] I numeri primi maggiori di 2 hanno successore pari;
    \item[(c)] Un numero primo è divisibile solo da 1 e da sé stesso;
    \item[(d)] I numeri primi formano un insieme discreto.
\end{enumerate}

\textbf{Svolgimento:}
Definiamo innanzitutto l'abbreviazione $NP(x) := N(x) \land Pr(x)$.

\begin{itemize}
    \item[(a)] \textbf{I numeri primi non formano un insieme denso.}
    Proprietà dell'ordine lineare:
    \[ \phi_{ANTIRIFL} := \forall x (NP(x) \to \neg(x < x)) \]
    \[ \phi_{TRICO} := \forall x \forall y ((NP(x) \land NP(y)) \to ((x < y) \lor (y < x) \lor (y = x))) \]
    \[ \phi_{TRANS} := \forall x \forall y \forall z ((NP(x) \land NP(y) \land NP(z) \land x < y \land y < z) \to (x < z)) \]
    \[ \phi_{LIN} := \phi_{ANTIRIFL} \land \phi_{TRICO} \land \phi_{TRANS} \]
    Proprietà di densità:
    \[ \phi' := \forall x \forall y ((NP(x) \land NP(y) \land (x < y)) \to \exists z \st (NP(z) \land (x < z) \land (z < y))) \]
    Soluzione:
    \[ \Phi_{NON\_DENSO} := \neg(\phi_{LIN} \land \phi') \]

    \item[(b)] \textbf{I numeri primi maggiori di 2 hanno successore pari.}
    \[ \forall x ((NP(x) \land (2 < x)) \to \exists s \st (N(s) \land Succ(s, x) \land Div(s, 2))) \]

    \item[(c)] \textbf{Un numero primo è divisibile solo da 1 e da sé stesso.}
    \[ \forall x (NP(x) \to \forall d (N(d) \land Div(x, d) \to ((d = 1) \lor (d = x)))) \]

    \item[(d)] \textbf{I numeri primi formano un insieme discreto.}
    \[ \phi'' := \forall x ((NP(x) \land \exists m \st (NP(m) \land (x < m))) \to \exists s \st (NP(s) \land (x < s) \land \forall t (NP(t) \land (x < t) \to \neg(t < s)))) \]
    Soluzione:
    \[ \phi_{DISC} := \phi_{LIN} \land \phi'' \]
\end{itemize}

\section*{Esercizio 2}
\textbf{Traccia:}
Si indichi, per ciascuna delle seguenti formule, se sono valide, insoddisfacibili o solo soddisfacibili. In ciascun caso, motivare semanticamente la risposta data.
Esibire, successivamente, una deduzione in Natural Deduction per ogni formula che risulta valida o per la sua negazione, se la formula risulta insoddisfacibile.
\begin{enumerate}
    \item[(a)] $(A \lor B) \to (\neg A \to B)$
    \item[(b)] $(A \to (B \lor \neg C)) \lor ((\neg A \land \neg B) \to \neg C)$
    \item[(c)] $\neg(A \land B) \land \neg(A \to \neg B)$
\end{enumerate}

\textbf{Svolgimento:}

\subsection*{(a) $(A \lor B) \to (\neg A \to B)$}
Semanticamente, la formula è equivalente a $(A \lor B) \to (A \lor B)$, che è una tautologia. Dunque è \textbf{Valida}.

\begin{prooftree}
  \AxiomC{\({[A \lor B]}_3\)}
  \AxiomC{\({[A]}_1\)}
  \AxiomC{\({[\neg A]}_2\)}
  \RuleLabel{\(\neg E\)}
  \BinaryInfC{\(\bot\)}
  \RuleLabel{\(\RAA\)}
  \UnaryInfC{\(B\)}
  \AxiomC{\({[B]}_1\)}
  \RuleLabel{\(\lor E_1\)}[\(A\), \(B\)]
  \TrinaryInfC{\(B\)}
  \RuleLabel{\(\to I_2\)}[\(\neg A\)]
  \UnaryInfC{\(\neg A \to B\)}
  \RuleLabel{\(\to I_3\)}[\(A \lor B\)]
  \UnaryInfC{\((A \lor B) \to (\neg A \to B)\)}
\end{prooftree}

\subsection*{(b) $(A \to (B \lor \neg C)) \lor ((\neg A \land \neg B) \to \neg C)$}
Sviluppando le implicazioni, otteniamo:
$(\neg A \lor B \lor \neg C) \lor ((A \lor B) \lor \neg C) \equiv (\neg A \lor A) \lor (B \lor B) \lor (\neg C \lor \neg C)$, che è banalmente vera in quanto contiene il terzo escluso $\neg A \lor A$. È \textbf{Valida}.

\begin{prooftree}
  \small
  \AxiomC{\({[\neg ((A \to (B \lor \neg C)) \lor ((\neg A \land \neg B) \to \neg C))]}_1\)}
  \AxiomC{\({[\neg A \land \neg B]}_3\)}
  \RuleLabel{\(\land E\)}
  \UnaryInfC{\(\neg A\)}
  \AxiomC{\({[A]}_4\)}
  \RuleLabel{\(\neg E\)}
  \BinaryInfC{\(\bot\)}
  \RuleLabel{\(\RAA_4\)}
  \UnaryInfC{\(B \lor \neg C\)}
  \RuleLabel{\(\to I\)}
  \UnaryInfC{\(A \to (B \lor \neg C)\)}
  \RuleLabel{\(\lor I\)}
  \UnaryInfC{\((A \to (B \lor \neg C)) \lor ((\neg A \land \neg B) \to \neg C)\)}
  \RuleLabel{\(\neg E\)}
  \BinaryInfC{\(\bot\)}
  \RuleLabel{\(\RAA_3\)}
  \UnaryInfC{\(\neg C\)}
  \RuleLabel{\(\to I\)}
  \UnaryInfC{\((\neg A \land \neg B) \to \neg C\)}
  \RuleLabel{\(\lor I\)}
  \UnaryInfC{\((A \to (B \lor \neg C)) \lor ((\neg A \land \neg B) \to \neg C)\)}
  \AxiomC{\({[\neg ((A \to (B \lor \neg C)) \lor ((\neg A \land \neg B) \to \neg C))]}_1\)}
  \RuleLabel{\(\neg E\)}
  \BinaryInfC{\(\bot\)}
  \RuleLabel{\(\RAA_1\)}
  \UnaryInfC{\((A \to (B \lor \neg C)) \lor ((\neg A \land \neg B) \to \neg C)\)}
\end{prooftree}

\subsection*{(c) $\neg(A \land B) \land \neg(A \to \neg B)$}
Sviluppando $\neg(A \to \neg B)$ si ottiene $\neg(\neg A \lor \neg B) \equiv A \land B$. La formula è quindi equivalente a $\neg(A \land B) \land (A \land B)$, che è una contraddizione. La formula è \textbf{Insoddisfacibile}.
Dimostriamo quindi la sua negazione $\neg(\neg(A \land B) \land \neg(A \to \neg B))$.

\begin{prooftree}
  \AxiomC{\({[\neg(A \land B) \land \neg(A \to \neg B)]}_1\)}
  \RuleLabel{\(\land E\)}
  \UnaryInfC{\(\neg(A \to \neg B)\)}
  \AxiomC{\({[A]}_3\)}
  \AxiomC{\({[B]}_2\)}
  \RuleLabel{\(\land I\)}
  \BinaryInfC{\(A \land B\)}
  \AxiomC{\({[\neg(A \land B) \land \neg(A \to \neg B)]}_1\)}
  \RuleLabel{\(\land E\)}
  \UnaryInfC{\(\neg(A \land B)\)}
  \RuleLabel{\(\neg E\)}
  \BinaryInfC{\(\bot\)}
  \RuleLabel{\(\to I_2\)}[\(B\)]
  \UnaryInfC{\(\neg B\)}
  \RuleLabel{\(\to I_3\)}[\(A\)]
  \UnaryInfC{\(A \to \neg B\)}
  \RuleLabel{\(\neg E\)}
  \BinaryInfC{\(\bot\)}
  \RuleLabel{\(\neg I_1\)}
  \UnaryInfC{\(\neg(\neg(A \land B) \land \neg(A \to \neg B))\)}
\end{prooftree}

\section*{Esercizio 4}
\textbf{Traccia:}
Si indichi, per ciascuna delle seguenti formule se sono valide, insoddisfacibili o solo soddisfacibili.
\begin{enumerate}
    \item[(a)] $\neg \exists x \st (A(x) \lor B(x)) \to \neg\neg \forall x \st \neg(A(x) \lor B(x))$
    \item[(b)] $(a=b) \to ((c=b) \to ((\forall x \st C(a,x)) \to (\exists x \st C(c,x))))$
\end{enumerate}

\textbf{Svolgimento:}

\subsection*{(a) $\neg \exists x \st (A(x) \lor B(x)) \to \neg\neg \forall x \st \neg(A(x) \lor B(x))$}
La formula è \textbf{Valida}. 

\begin{prooftree}
  \AxiomC{\({[\neg \forall x \st \neg(A(x) \lor B(x))]}_1\)}
  \AxiomC{\({[\neg \exists x \st (A(x) \lor B(x))]}_2\)}
  \AxiomC{\({[A(y) \lor B(y)]}_3\)}
  \RuleLabel{\(\exists I\)}
  \UnaryInfC{\(\exists x \st (A(x) \lor B(x))\)}
  \RuleLabel{\(\neg E\)}
  \BinaryInfC{\(\bot\)}
  \RuleLabel{\(\neg I_3\)}
  \UnaryInfC{\(\neg(A(y) \lor B(y))\)}
  \RuleLabel{\(\forall I\)}
  \UnaryInfC{\(\forall x \st \neg(A(x) \lor B(x))\)}
  \RuleLabel{\(\neg E\)}
  \BinaryInfC{\(\bot\)}
  \RuleLabel{\(\neg I_1\)}
  \UnaryInfC{\(\neg\neg \forall x \st \neg(A(x) \lor B(x))\)}
  \RuleLabel{\(\to I_2\)}
  \UnaryInfC{\(\neg \exists x \st (A(x) \lor B(x)) \to \neg\neg \forall x \st \neg(A(x) \lor B(x))\)}
\end{prooftree}

\subsection*{(b) $(a=b) \to ((c=b) \to ((\forall x \st C(a,x)) \to (\exists x \st C(c,x))))$}
Da $a=b$ e $c=b$ ricaviamo per simmetria e transitività $a=c$. Quindi $\forall x \st C(a,x)$ implica $\forall x \st C(c,x)$, da cui si deduce $\exists x \st C(c,x)$. La formula è \textbf{Valida}.

\begin{prooftree}
  \small
  \AxiomC{\({[a=b]}_3\)}
  \AxiomC{\({[c=b]}_2\)}
  \RuleLabel{\(Symm\)}
  \UnaryInfC{\(b=c\)}
  \RuleLabel{\(Trans\)}
  \BinaryInfC{\(a=c\)}
  \AxiomC{\({[\forall x \st C(a,x)]}_1\)}
  \RuleLabel{\(\forall E\)}
  \UnaryInfC{\(C(a,c)\)}
  \RuleLabel{\(Sub_\phi\)}
  \BinaryInfC{\(C(c,c)\)}
  \RuleLabel{\(\exists I\)}
  \UnaryInfC{\(\exists x \st C(c,x)\)}
  \RuleLabel{\(\to I_1\)}
  \UnaryInfC{\(\forall x \st C(a,x) \to \exists x \st C(c,x)\)}
  \RuleLabel{\(\to I_2\)}
  \UnaryInfC{\((c=b) \to (\forall x \st C(a,x) \to \exists x \st C(c,x))\)}
  \RuleLabel{\(\to I_3\)}
  \UnaryInfC{\((a=b) \to ((c=b) \to ((\forall x \st C(a,x)) \to (\exists x \st C(c,x))))\)}
\end{prooftree}
